\LARGE
\newpage
\subsection{Занятие 18. 2024.11.14}
Новые финали (Последний набор):
\begin{enumerate}
	\item ua (wa)
	\item uai (wai)
	\item u(e)i (wei)
	\item u(e)n (wen)\\
	Если 3-4 слоги сочетаются с инициалью, "e" редуцируется.
	\begin{example} \hfill\break
		问 - wèn\\
		王 - wáng
	\end{example}
	\item uan (wan)
	\item uang (wang)
	\item ueng (weng)
\end{enumerate}
\begin{newwords}\hfill
	\begin{enumerate}
		\item 会 - huì - уметь
		\begin{example}\hfill\break
			他会说汉语。\\
			他会说汉语吗？
		\end{example}
		\item 日语 - rìyǔ - японский язык
		\item 日本 - rìběn - Япония
		\item 英语 - yīngyǔ - английский язык
		\item 英国 - yīngguó - Англия
		\item 母语 - mǔyǔ - родной язык
		\item 名字 - míngzi - имя, кличка, название
		\begin{example}\hfill\break
			你叫什么名字？\\
			我叫吉洪。 / 我的名字叫吉洪。
		\end{example}
		\item 去 - qù - направляться (уходить/уезжать)
		\begin{example}\hfill\break
			我要去图书馆。 - Я хочу пойти в библиотеку.\\
			(图书馆 - túshūguǎn - библиотека)
		\end{example}
		\item 来 - lái - приехать, прийти, прибыть
		\item 来自 - láizì - приехать из ...
		\begin{example}\hfill\break
			他来自中国。
		\end{example}
		\item 俄罗斯 - éluósī - Россия (русский)
	\end{enumerate}
\end{newwords}